\addvspace{1cm}
\subsection{Rücktransformation, Inverse z-transformation}

\textbf{Inverse z-Transformation}

Auch bei der inversen Z-Transformation handelt es sich um eine lineare Transformation. Durch die Anwendung erhält man aus einer Bildfunktion $X(z)$ die zugehörige Urbildfolge $x[k]$ zurück. Die inverse z-Transformation ist definiert als:

\begin{equation}
    \label{eq:inverse_z_transformation}
    \mathcal{Z}^{-1}\{X(z)\} = x[k] = \frac{1}{2\pi j} \oint_{C} X(z) \cdot z^{k-1} dz
\end{equation}

Die Rücktransformation erfolgt über die Integration entlang eines geschlossenen Pfades $C$ im mathematisch positiven Sinn um den Ursprung im Konvergenzbereich der von $X(z)$. Das Ringintegral lässt sich zum Beispiel über das Residuenkalkül lösen, welches ein Werkzeug aus der Funktionentheorie ist. Diese Methode ist allerdings sehr kompliziert. Es gibt in der Regel einfachere Methoden, mit denen wir uns hier befassen wollen \cite[S.246]{book2}.
\addvspace{1cm}

\textbf{weitere Methoden zur Rücktransformation}

In der Praxis findet die Rücktransformation meist über \textbf{Korrespondenzen}, wie in der Tabelle \ref{tab:ZT_Korrespondenzen} angegeben statt.
Liegt $X(z)$ noch nicht in der Form der Korrespondenzen vor, so muss sie entweder über die in Abschnitt \ref{sec:Sätze_der_ZT} genannten Sätze umgeformt werden oder durch eine \textbf{Partialbruchzerlegung} in die benötigte Form gebracht werden \cite[S.246]{book2}.

Über die Definition der Potenzreihe lässt sich auch noch ein weiterer Ansatz zur Rücktransformation formulieren. Dabei kann das k-te Element der Urbildfolge $x[k]$ über die k-te Ableitung der Bildfunktion an der Stelle $\frac{1}{z}$ bestimmt werden:

Durch die Definition der z-Transformation ergibt sich:
\begin{equation}
    X(\frac{1}{z}) = \sum_{k=0}^{\infty} x[k] \cdot \left(\frac{1}{z}\right)^{-k} = \sum_{k=0}^{\infty} x[k] \cdot z^{k} = x[0] + x[1] \cdot z + x[2] \cdot z^2 + \ldots
\end{equation}

Leitet man die Funktion $X(\frac{1}{z})$ nun k-mal ab und bildet den Grenzwert für $z \rightarrow 0$, so erhält man $k! \cdot x[k]$:
\begin{align}
    \notag
    \lim_{z \to 0} F(\frac{1}{z}) &= x[0] + x[1] \cdot z + x[2] \cdot z^2 + x[3] \cdot z^3 + x[4] \cdot z^4 + \ldots = 0! x[0] \\
    \notag
    \lim_{z \to 0} \frac{d}{dz} F(\frac{1}{z}) &= 1 \cdot x[1] + 2 \cdot x[2] \cdot z + 3 \cdot x[3] \cdot z^2 + 4 \cdot x[4] \cdot z^3 + \ldots = 1! \cdot x[1] \\
    \notag
    \lim_{z \to 0} \frac{d^2}{dz^2} F(\frac{1}{z}) &= 2 \cdot 1 \cdot x[2] + 3 \cdot 2 \cdot x[3] \cdot z + 4 \cdot 3 \cdot x[4] \cdot z^2 + \ldots = 2! \cdot x[2]
     \\ \notag
    \lim_{z \to 0} \frac{d^3}{dz^3} F(\frac{1}{z}) &= 3 \cdot 2 \cdot 1 \cdot x[3] + 4 \cdot 3 \cdot 2 \cdot x[4] \cdot z + \ldots = 3! \cdot x[3] \\
    \notag
\end{align}

Über diesen Zusammenhang läst sich die folgende Berechnungsvorschrift für die Elemente der Urbildfolge $x[k]$ formulieren:
\begin{equation}
    \boxed{
    x[k] = \frac{1}{k!} \cdot \lim_{z \to 0} \frac{d^k}{dz^k} F(\frac{1}{z})
    }
\end{equation}  

Dazu sei noch erwähnt, dass diese Ableitungsmethode nur sinnvoll ist, wenn man ein Bildungsgesetz für die k-te Ableitung von $F(\frac{1}{z})$ finden kann.

\addvspace{1cm}

\textbf{Beispielberechnung zur inversen z-Transformation}

Im Folgenden betrachten wir die Rücktransformation der Bildfunktion $X(z) = \frac{2z+5}{z^2+5z-6}$, die für $|z| > 6$ konvergiert.
Die Berechnung der Polstellen, also der Nullstellen des Nennerpolynoms erfolgt bspw. über die pq-Formel:
\begin{equation}
    z = -\frac{5}{2} \pm \frac{7}{2} \Rightarrow  z = 1 \vee  z = -6
\end{equation}

Dadurch lässt sich die Funktion in Partialbrüche zerlegen. Da das Polyom einigermaßen harmlos ist, lassen sich die Werte für die Konstanten durch Einsetzen der Polstellen leicht berechnen:
\begin{equation}
\begin{aligned}
    &X(z) = \frac{2z+5}{z^2+5z-6} = \frac{A}{z-1} + \frac{B}{z+6} \Rightarrow 2z+5 = A(z+6) + B(z-1) \\
    &A = 1 \quad B = 1 \Rightarrow X(z) = \frac{1}{z-1} + \frac{1}{z+6} = \frac{1}{z} \cdot \left(\frac{z}{z-1} + \frac{z}{z-6}\right)
\end{aligned}
\end{equation}

An dieser Schreibweise kann man nun sofort die Korrespondenzen der Sprungfolge und der Exponentialfolge aus Tabelle \ref{tab:ZT_Korrespondenzen} ablesen, die nach dem Verschiebungssatz aus Abschnitt \ref{sec:Verschiebungssatz} durch den Vorfaktor $\frac{1}{z}$ um einen Abtastwert nach rechts verschoben sind:
\begin{equation}
\begin{aligned}
    &\mathcal{Z}^{-1}\left\{\frac{1}{z} \cdot \left(\frac{z}{z-1} + \frac{z}{z-6}\right)\right\} = z^{-1}\cdot\left( \epsilon[k] + \epsilon[k]\cdot(-6)^{k}\right) \\
    \\
    & x[k] = \epsilon[k-1] + \epsilon[k-1]\cdot(-6)^{k-1} \\ \\
    &\Rightarrow x[k] = \begin{cases}
        0 & k < 1 \\
        1 + (-6)^{k-1} & k \geq 1
    \end{cases}
\end{aligned}
\end{equation}

